\section{Context and Motivation}

Software systems are now built as large, heterogeneous, and continuously
changing artifacts. A single product may combine several programming languages,
frameworks, generated files, configuration formats, build scripts, tests,
front-end components, backend services, data-access layers, and infrastructure
definitions. This heterogeneity is not merely an implementation detail. It
changes how developers understand systems, how maintainers reason about
evolution, and how teams assess whether a codebase still follows expected
architectural and organizational conventions.

Program comprehension has therefore become one of the central activities in
software maintenance. Developers spend substantial time locating relevant
artifacts, recognizing responsibilities, understanding dependencies, and
deciding whether a given file conforms to project-specific expectations.
Traditional static analysis tools support this process by extracting metrics,
detecting code smells, checking rules, and reporting violations. Software
visualization tools complement these mechanisms by exposing dependencies,
architectural structures, evolution patterns, or fault-localization evidence in
forms that are easier to inspect than raw source code
\cite{pinzger2008tool,novais2013sourceminer,kobayashi2013sarf,qayum2022finecodeanalyzer,sun2024sflvis}.

At the same time, machine learning has expanded the range of possible source
code analyses. Modern models can process tokens, abstract syntax trees, graphs,
embeddings, and long code sequences for tasks such as code classification,
defect prediction, clone detection, summarization, and vulnerability analysis
\cite{allamanis2018survey,alon2019code2vec,feng2020codebert,guo2021graphcodebert,yang2025sparsecoder}.
These representations are powerful, but they often depend on language-specific
parsers, tokenizers, semantic analyzers, or pretraining corpora. Such
requirements can make the first layer of analysis harder to reproduce across
multiple ecosystems and harder to integrate into lightweight tooling.

This qualification project investigates a complementary path: representing
source code as visual signatures. In particular, it studies \textit{source code
minimaps}, compact image-like representations generated from source files. A
minimap preserves spatial and structural cues such as indentation, line density,
punctuation, block delimiters, and layout regularities, while lexical content can
be anonymized. This idea is inspired by the minimap widgets that developers
already know from modern editors, but the purpose is different. In this research,
the minimap is not primarily a user-interface element for manual navigation. It
is a machine-readable visual abstraction that can be processed with texture
descriptors, computer vision models, deep learning architectures, and
explainable artificial intelligence techniques.

The motivation for this representation is twofold. First, software artifacts
exhibit visual regularities. A Java entity, a test class, a JSON configuration, a
large generated file, a shell script, and an SVG file may differ substantially in
their visual structure even before their full lexical content is interpreted.
Second, visual abstractions can provide a common input format across languages
and file types. This does not remove the value of parsers, static analyzers, or
semantic models. Rather, it provides a lightweight first representation that can
be used to classify, cluster, inspect, or prioritize artifacts before deeper
analysis is applied.

The project has already produced preliminary evidence in this direction. An
initial study used source code minimaps, Local Binary Patterns, and traditional
classifiers to detect software features, with Random Forest reaching accuracy
near 93\% in the evaluated setting \cite{gebara2025iwssip}. A second study
scaled the approach to 685,906 minimap images extracted from 100 open-source
projects across 10 programming languages, using deep learning models and
Integrated Gradients for interpretation \cite{gebara2025ciarp}. A tool-support
line, named \textit{MinimapSignatures}, then brought the approach into Visual
Studio Code and IntelliJ IDEA through IDE plugins connected to a Python backend
\cite{gebara2026minimapsignatures}. This qualification document consolidates
these results into a coherent doctoral proposal and defines the remaining path
toward the thesis.

\section{Problem Statement}

The general problem addressed in this research is the difficulty of performing
scalable, reproducible, and language-agnostic source code analysis across
heterogeneous software repositories. More specifically, the project focuses on
the following tension:

\begin{quote}
Software projects contain rich structural and stylistic signals that are useful
for classification, comprehension, quality assessment, and tool support, but
many existing representations either require language-specific processing or are
not designed to be compact, anonymized, and directly reusable by visual learning
models.
\end{quote}

Traditional static analysis is effective when the target language, framework,
and rule set are known. However, such tools typically require parsers,
language-specific rules, and expert configuration. Software visualization
approaches expose structures to human analysts, but many of them are not
designed as machine-readable representations for large-scale classification.
Machine learning representations for code, in turn, commonly rely on tokens,
abstract syntax trees, graphs, or pretrained language models; these
representations may require significant preprocessing and may retain lexical
content that is sensitive in industrial contexts.

The project therefore asks whether a visual representation of source files can
serve as an intermediate abstraction: simple enough to be generated across many
languages and file types, compact enough to support large datasets, structured
enough to retain useful information, and anonymizable enough to reduce exposure
of raw source code. Source code minimaps are proposed as this abstraction.

\section{Research Gap}

The literature contains several lines of work related to this proposal. Software
visualization has produced tools for dependency inspection, architecture
recovery, feature evolution, code-smell visualization, and fault localization
\cite{pinzger2008tool,novais2013sourceminer,kobayashi2013sarf,reis2022code,sun2024sflvis}.
Machine learning for source code has explored token sequences, graph
representations, embeddings, and transformer-based models
\cite{allamanis2018survey,feng2020codebert,guo2021graphcodebert,yang2025sparsecoder}.
More recent work has also begun to treat source code as visual data, as in
CV4Code \cite{shi2023cv4code}, and language-agnostic visual inspection tools
such as CodePanorama \cite{etter2022codepanorama} have shown the relevance of
zoomed-out representations.

Despite these advances, four gaps remain relevant for this doctoral project.
First, many visual tools are primarily designed for human inspection rather than
for systematic machine learning over large corpora. Second, many learning-based
representations depend on language-specific preprocessing or preserve lexical
content that may be sensitive. Third, there is still limited empirical evidence
about how much structural information compact minimap images preserve across
many projects, languages, and prediction objectives. Fourth, minimap-based
analysis has rarely been operationalized in the environments where developers
actually work, such as integrated development environments.

This project positions source code minimaps at the intersection of these gaps.
It treats minimaps as compact visual signatures, evaluates them empirically with
traditional and deep learning models, interprets model behavior with explainable
AI, and integrates the resulting analysis into IDE-based tooling.

\section{Hypothesis}

The central hypothesis of this doctoral project is:

\begin{quote}
Anonymized source code minimaps preserve enough structural and visual
information to support scalable and language-agnostic software analysis, while
reducing dependence on language-specific parsers and expert-defined rules.
\end{quote}

The hypothesis has three important boundaries. First, it does not claim that
minimaps replace semantic analysis. Source code meaning still depends on
language constructs, dependencies, runtime behavior, and domain knowledge.
Second, it does not claim that anonymization provides a formal privacy guarantee.
The goal is to reduce lexical exposure while preserving structural cues; privacy
must still be evaluated carefully. Third, it does not claim that a single visual
representation is sufficient for all software-engineering tasks. Instead, it
claims that minimaps can provide a useful, reproducible, and extensible first
representation for a family of analysis tasks.

\section{Research Questions}

The project is organized around five research questions:

\begin{description}
    \item[RQ1.] To what extent can source code minimaps preserve discriminative
    structural cues for source artifact classification across programming
    languages?

    \item[RQ2.] How do handcrafted texture descriptors compare with learned
    visual representations for minimap-based software analysis?

    \item[RQ3.] Which visual regions or patterns do deep models use when
    classifying source code minimaps, and how can explainable AI support the
    interpretation of these decisions?

    \item[RQ4.] How does anonymization affect the balance between
    confidentiality and predictive performance in minimap-based analysis?

    \item[RQ5.] How can source code minimap analysis be integrated into practical
    IDE workflows for developer-facing feedback?
\end{description}

These questions connect representation, empirical evaluation, interpretation,
privacy, and tool support. RQ1 and RQ2 address the feasibility and modeling
capacity of minimaps. RQ3 addresses transparency. RQ4 addresses the ethical and
practical need to reduce exposure of source code. RQ5 addresses the transition
from offline experiments to software-engineering practice.

\section{Objectives}

The general objective of this doctoral project is to define, evaluate, and
operationalize a source code minimap approach for visual software analysis.

The specific objectives are:

\begin{enumerate}
    \item Define a reproducible method for generating source code minimaps from
    heterogeneous software repositories.

    \item Design and evaluate anonymization strategies that suppress lexical
    content while preserving structural and visual cues.

    \item Build datasets and catalogs that support supervised learning tasks
    over minimap images.

    \item Compare handcrafted texture descriptors and deep visual models for
    minimap-based classification.

    \item Investigate explainable AI techniques to identify which visual regions
    contribute to model predictions.

    \item Develop tool support that integrates minimap generation and inference
    into developer environments.

    \item Extend the current evaluation toward software quality, code-smell
    detection, architectural deviation analysis, and project-specific model
    adaptation.
\end{enumerate}

\section{Expected Contributions}

The expected contributions of the thesis are both scientific and practical.
Scientifically, the project aims to establish source code minimaps as a
reproducible visual representation for software analysis, supported by empirical
evidence across datasets, models, and tasks. It also contributes to the growing
body of work on visual representations of code by studying anonymization,
class imbalance, cross-language applicability, and explainability.

Practically, the project aims to produce datasets, scripts, trained-model
pipelines, and IDE-integrated tools that can be reused and extended by other
researchers and practitioners. The MinimapSignatures ecosystem is particularly
important in this regard because it moves the approach from offline notebooks
and experiments into development environments.

Table~\ref{tab:intro-contributions} summarizes the main contribution areas and
their current status.

\begin{table}[ht]
\centering
\caption{Expected contribution areas and current status.}
\label{tab:intro-contributions}
\begin{tabular}{p{0.30\textwidth}p{0.38\textwidth}p{0.20\textwidth}}
\toprule
\textbf{Contribution area} & \textbf{Description} & \textbf{Status} \\
\midrule
Representation & Source code minimaps as compact visual signatures of software artifacts. & Validated (two studies) \\
Anonymization & Character-to-pixel mapping that suppresses lexical content while preserving layout signals. & Validated; ablation study planned \\
Datasets & 685,906-image corpus from 100 projects across 10 languages; archived on Zenodo. & Available \\
Modeling & LBP + shallow classifiers (Study 1); ResNet-18 and EfficientNet-B0 (Study 2); custom multitask CNN (Stage 3). & Validated \\
Explainability & Integrated Gradients attribution maps for Type, Project, and Author objectives. & Initial results \\
Tool support & VS Code and IntelliJ IDEA plugins connected to a FastAPI/PyTorch backend; functional prototype. & Prototype validated \\
Quality tasks & Extension toward code smells, quality indicators, and architectural deviations. & Planned \\
\bottomrule
\end{tabular}
\fautor
\end{table}

\section{Scope and Assumptions}

The scope of this qualification proposal is intentionally focused on file-level
visual signatures. A source code minimap represents one artifact at a time. It
does not directly encode the full dependency graph of a project, the runtime
behavior of a system, the semantics of API calls, or the history of a design
decision. These aspects remain important, but they belong to complementary
analysis layers. The initial research question is whether a compact visual
projection of a single artifact already contains useful signals for software
analysis.

This scope has practical advantages. File-level artifacts are the units that
developers often inspect during maintenance and review. They are also convenient
for dataset generation because repositories can be traversed automatically and
each file can be rendered independently. Moreover, file-level analysis is a
reasonable first step for IDE integration. A developer can trigger a command on
the active editor and receive immediate feedback without submitting an entire
repository.

The project also assumes that visual structure is meaningful even when lexical
content is suppressed. This assumption is plausible because software artifacts
are not random text. Formatting conventions, indentation, block delimiters,
annotation placement, repeated declarations, template regions, and generated
sections produce visual patterns. Nevertheless, the assumption must be tested
empirically. If minimaps only worked when lexical content was preserved, the
privacy and generality motivation would be weak. For this reason, anonymization
is treated as part of the representation, not as an optional post-processing
step.

Another assumption is that a language-light first representation is useful even
when language-specific analysis remains necessary for deeper tasks. A minimap
may identify that a file visually resembles a test class, but it cannot prove
that the tests assert meaningful behavior. It may show that a source artifact
deviates from common project patterns, but it cannot alone determine whether the
deviation is a defect or an intentional design choice. The intended use is
therefore decision support: minimaps can help prioritize, classify, cluster, and
explain artifacts, while human judgment and complementary tools remain central.

Finally, the project assumes that practical adoption depends on tool support.
An offline dataset and model can demonstrate scientific feasibility, but a
software-engineering technique becomes more meaningful when it can be connected
to actual workflows. MinimapSignatures is included in the scope of the doctorate
for this reason. It provides a path from representation learning to developer
interaction, and it creates a platform for future studies about usefulness,
trust, and explanation design.

\section{Document Organization}

The remainder of this qualification document is organized as follows.
Chapter~\ref{chapter:foundations} presents the theoretical foundations required
to understand the proposal, including software visualization, static analysis,
source code representations, machine learning, computer vision, and explainable
AI. Chapter~\ref{chapter:related-work} reviews related work and synthesizes the
literature gaps that motivate the research. Chapter~\ref{chapter:methodology}
describes the proposed methodology for generating, anonymizing, cataloging, and
analyzing source code minimaps. Chapter~\ref{chapter:preliminary-results}
consolidates preliminary studies and results already produced during the
doctorate. Chapter~\ref{chapter:tool-support} presents MinimapSignatures, the
tool ecosystem that integrates the approach into IDE workflows.
Chapter~\ref{chapter:research-plan} defines the remaining research plan, risks,
threats to validity, and ethical considerations. Finally,
Chapter~\ref{chapter:final-remarks} summarizes the proposal and expected thesis
contributions.
